%% ma: L12-B08-0150
%% lop: 12
%% chuong: 2
%% bai: 8
%% dang: 12.8.50
%% loai: TN4
%% muc: VD
%% dap_an: "C"
%% nguon: "CĐ2.3.pdf giáo viên gửi 27/10/2026 (M07, Lê Minh Khôi – Nguyễn Anh Khôi, Chương II Chủ đề 3 Biểu thức tọa độ của các phép toán vectơ), Câu 66"
%% kien_thuc: "Bài 8 (điểm chia đoạn thẳng, độ dài), lớp 10 (tính chất đường phân giác)"
%% trang_thai: chinh-thuc
%% ghi_chu: "Giáo viên đã duyệt 27/10/2026"
\begin{ex}
Trong không gian $Oxyz$, cho ba điểm $A(-1;2;2)$, $B(0;4;0)$, $C(-1;1;2)$. Gọi $D(a;b;c)$ là chân đường phân giác trong của tam giác $ABC$ kẻ từ $A$. Tính giá trị của biểu thức $P=a+b+c$.
\choice{$P=\dfrac32$}{$P=\dfrac72$}{\True $P=\dfrac52$}{$P=\dfrac92$}
\loigiai{$AB=3$, $AC=1$ nên $\overrightarrow{DB}=-3\overrightarrow{DC}\Rightarrow D=\dfrac{B+3C}4=\left(-\dfrac34;\dfrac74;\dfrac32\right)$, $P=\dfrac52$.}
\end{ex}
